\section*{Checklist}

\begin{enumerate}

\item For all authors...
\begin{enumerate}
  \item Do the main claims made in the abstract and introduction accurately reflect the paper's contributions and scope?
    \answerYes{}
  \item Did you describe the limitations of your work?
    \answerYes{}, please refer to the description in the supplementary material or the following content.

    \paragraph{Dataset.}
    Due to the high cost of collecting and annotating data, there are currently only 83 TV dramas in \TVD.
    As analysed in \cref{sec:experiment_res}, the genre distribution of \TVD is unbalanced.
    The text data of \TVD (i.e., the plot text and Bullet Screen Comments) is in Chinese only, and heavily reflects Chinese culture and expression habits. Therefore, \TVD may exhibit gaps between real-world data from other social backgrounds.
    Also, although we adhere to strict standards in our data collection pipeline, we are not able to manually check all the samples, so \TVD may contain a small percentage of annotation errors.

    \paragraph{Framework.}
    To reduce the environmental impact and to highlight scalability, the framework implemented in this study adopts the very popular techniques rather than the most advanced ones.
    As identified in \cref{sec:experiment_res}, our framework may aggregate existing distribution bias in the dataset, which should be treated very carefully.
    In the task of Plot Text Generation, we did not conduct manual evaluation as our budget has run out during the creation of \TVD (annotation expense, copyright overhead, etc.).
    However, we should remind readers that human judgment is still needed in some cases.

    \paragraph{Experiments.}
    Our framework, together with \TVD, could be used to explore a wide range of tasks, but due to time limit we only tested three most representative ones.
    We made several counter-intuition observations in  \cref{sec:experiment_res}. Although we have provided reasonable guesses on their reasons, we admit a very in-depth investigation still requires a considerable amount of work.

  \item Did you discuss any potential negative societal impacts of your work?
    \answerYes{}, please refer to the description in the supplementary material or the following content.
    
    Most of the TV dramas included in \TVD are produced in a fictional context, and the relevant content may have the specific bias introduced by the original creators.
    Meanwhile, the filming techniques and editing strategies commonly used in film and television photography will also highlight some dramatic plots and may have a preference for some content.
    In \TVD,  the Bullet Screen Comments carry the opinions of the user, which may therefore contain personal bias.
    The above bias may be inherited or even aggregated by the downstream algorithms.
    Thus, we urge future users of \TVD to pay special  attention to this matter.
  \item Have you read the ethics review guidelines and ensured that your paper conforms to them?
    \answerYes{}, please refer to the description in the supplementary material or the following content.
    
    Our data annotation and release plan has been examined and approved by relevant institutional committees.
    All annotators participating the creation of \TVD have agreed relevant terms and received a set of pre-annotation training.
    They were provided with remuneration that is significantly higher than the local minimum wage.
    At different stages of data collection, we complied with corresponding legal and ethical requirements.
    In order to protect the privacy of TV drama viewer, we have also anonymised all data (especially Bullet Screen Comments), ensuring no information that can uniquely identify users is excluded from our dataset and research process.
    We have obtained all data copyrights relevant to this paper.
    We have prepared detailed instructions, licenses, and data use agreement for future users of \TVD, aiming to minimise the possibility of any kind of misuse (e.g., to train a malicious Bullet Screen Comment writer).
\end{enumerate}

\item If you are including theoretical results...
\begin{enumerate}
  \item Did you state the full set of assumptions of all theoretical results?
    \answerYes{}
	\item Did you include complete proofs of all theoretical results?
    \answerYes{}
\end{enumerate}

\item If you ran experiments (e.g. for benchmarks)...
\begin{enumerate}
  \item Did you include the code, data, and instructions needed to reproduce the main experimental results (either in the supplemental material or as a URL)?
    \answerYes{}, please refer to \url{https://ptvd.github.io/}.
  \item Did you specify all the training details (e.g., data splits, hyperparameters, how they were chosen)?
    \answerYes{}, please refer to the description in the supplementary material or the following content. In the three experiments, we used vision and text encoders are vit\_base\_p16\_224\footnote{\url{https://huggingface.co/google/vit-base-patch16-224}} and bert-base-chinese\footnote{\url{https://huggingface.co/bert-base-chinese}} respectively.
    In the image-text pre-training model, we use the AdamW optimizer~\citep{adamw2017axriv} with the learning rate set to 1e-5 for the pre-training model.
    The warm-up ratio is set to 10\%, and the learning rate decays linearly to 0 after 10\% of the total training steps.
    The batch size, hidden size, and number of heads are set to 4096, 768, and 12, respectively.
    We pre-trained our model for 100K steps on 8 NVIDIA A100 GPUs, which takes around 5 days.
    In the video-text pre-training model, since the video representation is obtained by inputting four frames of the video into the image encoder in our framework, we have completely followed the settings of all image-text pre-training models. 
    This model is also pre-trained for 100K steps on 8 NVIDIA A100 GPUs, which takes about 8 days. 
	\item Did you report error bars (e.g., with respect to the random seed after running experiments multiple times)?
    \answerNo{}, Because the training data size of METER is very large, the choice of a random seed usually does not have a large impact. At the same time, pretraining of large model often requires extremely expensive resources, so we do not report error values and release the model weights~\citep{meter2022cvpr}.
	\item Did you include the total amount of compute and the type of resources used (e.g., type of GPUs, internal cluster, or cloud provider)?
    \answerYes{}, see above.
\end{enumerate}

\item If you are using existing assets (e.g., code, data, models) or curating/releasing new assets...
\begin{enumerate}
  \item If your work uses existing assets, did you cite the creators?
    \answerYes{}
  \item Did you mention the license of the assets?
    \answerYes{}
  \item Did you include any new assets either in the supplemental material or as a URL?
    \answerYes{}
  \item Did you discuss whether and how consent was obtained from people whose data you're using/curating?
    \answerYes{}
  \item Did you discuss whether the data you are using/curating contains personally identifiable information or offensive content?
    \answerYes{}
\end{enumerate}

\item If you used crowdsourcing or conducted research with human subjects...
\begin{enumerate}
  \item Did you include the full text of instructions given to participants and screenshots, if applicable?
    \answerYes{}, please refer to \cref{sec:construction}.
  \item Did you describe any potential participant risks, with links to Institutional Review Board (IRB) approvals, if applicable?
    \answerYes{}
  \item Did you include the estimated hourly wage paid to participants and the total amount spent on participant compensation?
    \answerYes{}, We implemented the annotation task through a web crowdsourcing platform\footnote{\url{https://soho.qq.com/}.}, which cost US \$3,770 in total, or US \$0.17 per treaty on average. If each annotation takes 1-2 minutes to calculate, the hourly rate of annotation is about \$6.80.
\end{enumerate}

\end{enumerate}